\section{Supervision and Fault Tolerance}
\label{sec:Supervision}

Durga runtime components communicate through Kafka topics, which act as
durable mailboxes. This architecture enables an Erlang/OTP--like supervision
model: workers can crash and restart independently while Kafka preserves
unconsumed records according to topic durability and retention settings.

\subsection{Architecture}

\begin{itemize}
\item Every worker consumes from a dedicated Kafka topic (its inbox) and
      produces to downstream topics.
\item Kafka persists messages for a configurable retention period. A crashed
      worker's unconsumed messages remain in its topic.
\item Consumer group rebalancing shifts partitions to surviving instances when
      a worker pod exits.
\item The scaffolder's \texttt{--separate-workers} flag generates one
      Kubernetes Deployment per BPMN task, giving each worker its own
      fault domain.
\end{itemize}

\subsection{Per-worker isolation}

When \texttt{--separate-workers} is enabled:

\begin{verbatim}
./durga durga-tools/src/test/resources/bpmn/invoice_receipt.bpmn \
--separate-workers
\end{verbatim}

Each generated project includes:

\begin{itemize}
\item \texttt{RegisterInvoiceStandalone.java}, \texttt{ReviewInvoiceStandalone.java}, \\
      \texttt{NotifyRequesterStandalone.java} --- one standalone worker per BPMN
      task, each with its own Kafka consumer/producer and poll loop.
\item \texttt{InvoiceReceiptWorkerBootstrap.java} --- entry point that reads
      the \texttt{WORKER\_ID} environment variable and starts the matching
      worker.
\item \texttt{Dockerfile.<task>} --- per-worker Docker image building from
      the same JAR, with \texttt{WORKER\_ID} pre-set.
\item \texttt{k8s.<task>.yml} --- per-worker Kubernetes Deployment with
      isolated resource limits, health probes, and \texttt{WORKER\_ID}
      environment variable.
\end{itemize}

\subsection{Supervision policies}

\begin{longtable}{|>{\raggedright\arraybackslash}p{0.22\textwidth}|>{\raggedright\arraybackslash}p{0.70\textwidth}|}
\hline
\rowcolor[gray]{0.9}\textbf{Policy} & \textbf{Implementation} \\
\hline
Crash recovery & Kafka consumer offset tracking. Worker restarts replay from
last committed offset. \\
\hline
Fault isolation & Per-worker Deployments. One worker's OOM does not affect
others. \\
\hline
Health signals & Exec-based liveness/readiness probes per pod (check process
alive). Extendable to application-level health via
\texttt{quarkus-smallrye-health}. \\
\hline
Scaling & \texttt{kubectl scale} per Deployment. Bounded by topic partition
count. \\
\hline
Dead letter & Plugin executor routes failures to \texttt{\%task\%\_dlq}
topic. Standalone workers log errors and continue. \\
\hline
Backpressure & Kafka consumer group throttles ingestion. Consumer lag
available via JMX metrics. \\
\hline
\end{longtable}

\subsection{Recovery guarantees}

\begin{description}
\item[Message durability] Messages are not lost on worker crash. Kafka
      retains uncommitted records according to the topic's replication,
      acknowledgement and retention configuration.
\item[At-most-once vs at-least-once] Standalone workers commit after
      processing (\texttt{commitSync}). The poll loop catches and logs
      errors without crashing. Plugin-generated executors route failures
      to the DLQ. The architecture provides at-least-once delivery with
      idempotent output and external side-effect handling being the application
      owner's responsibility.
\item[Process continuity] If the starter crashes, the process never starts
      (one-shot job, retried by K8s). If a worker crashes, Kafka consumer
      group rebalancing reassigns its partitions. If the final orchestrator
      crashes, the process remains in-flight until a new orchestrator pod
      picks up.
\end{description}

\subsection{Limitations}

\begin{itemize}
\item Worker state is external (Kafka). In-memory state such as running
      aggregations or cancellation registrations is lost on crash unless that
      specific generated handler persists it.
\item The \texttt{--separate-workers} model generates poll-loop workers
      rather than Quarkus CDI beans. These trade CDI features
      (dependency injection, reactive messaging) for isolated fault
      domains and smaller per-pod memory footprints.
\item Cross-worker coordination (AND joins, subprocess completion) relies
      on the \texttt{ProcessStateStore}, which is itself a Kafka
      consumer. Its crash characteristics are the same as any other
      worker.
\item Kafka durability does not imply business-level exactly-once semantics.
      Handlers that call databases, HTTP services or file systems must be
      idempotent or protected by application-level deduplication.
\end{itemize}
